\chapter{Estrutura do Repositório e Instruções de Execução}
\label{anexo1}

\section{Estrutura de Ficheiros}
\label{sec:estrutura_ficheiros}

O repositório está organizado da seguinte forma:

\begin{verbatim}
ai-killer-sudoku/
  src/
    prolog/
      sudoku.pl         - Predicados base do Sudoku classico
      dfid.pl           - Solucionador DFID com MRV
      astar.pl          - Solucionador A* (h=celulas vazias)
      killer_sudoku.pl  - Extensao Killer Sudoku (DFID + A*)
    python/
      board.py          - Representacao da grelha e utilitarios
      puzzles.py        - Caminhos dos puzzles
      sa_solver.py      - Solucionador SA
      ga_solver.py      - Solucionador GA
  data/
    input/
      classic_easy_01.txt
      classic_medium_01.txt
      classic_hard_01.txt
      classic_expert_01.txt
      killer_easy_01.txt
    output/
      benchmark_results.txt
      sa_results.txt
      ga_results.txt
  scripts/
    benchmark.sh        - Script de benchmark completo
\end{verbatim}

\section{Instruções de Execução}
\label{sec:instrucoes}

\paragraph{Requisitos.}
SWI-Prolog (\texttt{swipl}) e Python 3.10+ devem estar instalados e disponíveis no \texttt{PATH}. Não são necessárias bibliotecas Python externas.

\paragraph{Benchmark Completo.}
Para executar todos os solucionadores em todos os puzzles e gerar o relatório de resultados:
\begin{verbatim}
  bash scripts/benchmark.sh
\end{verbatim}

\paragraph{Solucionadores Prolog individualmente.}
\begin{verbatim}
  # DFID — puzzle fácil
  swipl -q -s src/prolog/dfid.pl \
        -g "solve_file('data/input/classic_easy_01.txt', _), halt"

  # A* — puzzle especialista
  swipl -q -s src/prolog/astar.pl \
        -g "solve_file('data/input/classic_expert_01.txt', _), halt"

  # Killer Sudoku (DFID + A*)
  swipl -q -s src/prolog/killer_sudoku.pl \
        -g "solve_killer_file(easy_01, _), halt"
\end{verbatim}

\paragraph{Solucionadores Python individualmente.}
\begin{verbatim}
  # SA — puzzle médio, 5 execuções
  python3 src/python/sa_solver.py --puzzle medium --runs 5

  # GA — puzzle fácil, população 300, 5000 gerações
  python3 src/python/ga_solver.py --puzzle easy --pop 300 \
          --generations 5000 --runs 3

  # SA — Killer Sudoku
  python3 src/python/sa_solver.py --puzzle killer --runs 3

  # GA — Killer Sudoku
  python3 src/python/ga_solver.py --puzzle killer --pop 200 \
          --generations 10000 --runs 3
\end{verbatim}
